\section{Security Statements and Proof Outlines}\label{app:security}
This appendix states the claims of Section~\ref{sec:threat} and says, for
each, which part is a Rocq theorem and which part is an argument on
paper. Nothing in this appendix is mechanised beyond the theorems it
cites by name. For those, Rocq's \texttt{Print Assumptions} reports the
soundness, zero-knowledge, and end-to-end theorems closed under the
global context. Three completeness lemmas use functional extensionality,
and the 2048-bit instance uses Rocq's primitive integers to check its
primality certificates (\rocqfile{bench/assumptions.sh}).

\newcounter{secclaim}
\textbf{Model.} The group $\langle g \rangle$ has prime order $q$ and
DDH holds in it. The hash is a random oracle with range $\mathbb{Z}_q$,
with separate domains for the ballots of each credential, for the
certificate, and for the commitment key. (Our driver separates ballots
by a credential prefix and the three kinds of query by their formats
only, and it reduces a $256$-bit hash modulo the $256$-bit $q$, which at
worst doubles the bound below.) The election parameters ($g$,
$h$, $m$, the commitment key) are fixed before the adversary runs. The
board is public and append-only, and it accepts one ballot per
credential. The trustees share $x$ with $h = g^x$ as in
Section~\ref{sec:discussion}.

\textbf{Result function.} For rankings $s_1, \ldots, s_N$ let $D$ be
their count matrix. Let $w(D)$ be the first Schulze winner of $D$ in
candidate order (\rocq{winner_of}) and $\Pi(D)$ the witness paths that
\rocq{witness_path} computes, for every pair, from $D$. The result of
the election is $\mathcal{L}(D) = (w(D), \Pi(D))$. Public privacy is
privacy with respect to this result function, and internal privacy is
privacy with respect to $D$ itself, as in Kryvos~\cite{kryvos2022}.

\subsection{Verifiability}
\refstepcounter{secclaim}\label{prop:verif}
\textbf{Proposition~\thesecclaim{} (first ballot).} Let $g \ne 1$ and
$q > \max(4m + 2,\, 2N^2 + 3N)$. Let an adversary that may control every
voter and every trustee make at most $Q$ oracle queries and output a
board with $N$ ballots that \rocq{nizk_verify_ballot} accepts and a
certificate for $w$ that \rocq{nizk_verify_tally} accepts on their
aggregate. Then, except with probability at most $(Q + N + 1)/q$, there
are position vectors $s_1, \ldots, s_N \in [0, m]^m$ such that the bit
ciphertexts of ballot $k$ encrypt the bits $[s_{k,i} < s_{k,j}]$, and
$w$ is a Schulze winner of the count matrix of $s_1, \ldots, s_N$.

\emph{Proof.} Call the ciphertexts of a ballot \emph{valid} if they
encrypt the bits of some position vector in $[0,m]^m$. Validity depends
on the ciphertexts and on $x$ only. By \rocq{ballot_sound}, if the
ciphertexts are not valid, then for every choice of announcements at
most one challenge has responses that the verifier accepts: two such
challenges would give a position vector. The hypotheses of the theorem
hold because $g \ne 1$ has order $q$ (\rocq{pow_injective}) and $q >
4m + 2$. The non-interactive verifier accepts only if the challenge is
the oracle's value at the credential, the ciphertexts, and the
announcements. So a query on invalid ciphertexts returns the one
accepting challenge with probability at most $1/q$. If all $N$ ballots
are valid, their aggregate encrypts a count matrix $D$ with entries in
$[0, N]$ (\rocq{cdec_aggregate}). Call the pair of the aggregate and $w$
\emph{correct} if $w$ is a Schulze winner of $D$, again a property of
the board and $x$ only. By \rocq{verify_tally_sound} and
\rocq{cert_winner}, if the pair is not correct, then for every choice of
the certificate's ciphertexts, paths, and announcements at most one
challenge is accepted, and the verifier takes the challenge from the
oracle at exactly these values. The adversary makes $Q$ queries and the
verifier $N + 1$. A union bound over them gives the claim. \hfill$\Box$

The proposition counts the ballots that are on the board. Who may post a
ballot, and whether a ballot holds the ranking its voter intended, are
outside it (Section~\ref{sec:threat}). When several candidates are
Schulze winners, the proposition does not say which one is announced
(Section~\ref{sec:discussion}).

\textbf{Batched ballots.} The batched ballot is an argument: a
commitment has many openings, so validity cannot be read off the first
message, and the step to Fiat--Shamir needs an extractor that rewinds.
\rocq{ballot3_sound} gives the tree of transcripts from which the
extractor succeeds. At the challenge $\beta$ the tree is more specific
than it has to be. Fix the ciphertexts $c$ and a vector $b$, and say that
the \emph{link} holds at $\beta$ if $\prod_k c_k^{\beta_k} = (g^R,\,
g^{\langle \beta, b \rangle} h^R)$ for some $R$. Both sides are
homomorphic in $\beta$. So if the link holds at $\beta^{(1)}, \ldots,
\beta^{(t)}$, it holds at every linear combination of them, and if these
vectors span $\mathbb{Z}_q^{m^2}$, it holds at every unit vector, which
says that $c_k$ encrypts $b_k$. The theorem's vectors, a base vector and
its $m^2$ neighbours in one coordinate each, are one spanning family.
Any $m^2$ linearly independent challenge vectors would do, and a set of
vectors fails to span only if it lies in a proper subspace, which a
uniform $\beta$ hits with probability at most $1/q$. Together with the
thresholds at the other challenges ($m^2$ values of $y$, $L + 2$ of $z$,
two of $c$, $2(L+3)+1$ of $x$), this is special soundness for general
access structures, for which Attema et al.~\cite{afkr23} show that
Fiat--Shamir keeps knowledge soundness with a loss linear in the number
of queries. Three things remain to be done before this is a proof: to
restate \rocq{ballot3_sound} for linearly independent vectors, to check
the first round, which has three challenges from separate queries on the
same input, against the hypotheses of~\cite{afkr23}, and to derive the
bound. Until then Proposition~\ref{prop:verif} is proved for the first
ballot only, and for the batched ballot we rely on
\rocq{schulze_end_to_end3} in its rewinding form.

\subsection{Privacy}
\refstepcounter{secclaim}\label{prop:priv}
\textbf{Claim~\thesecclaim{} (privacy; proof outline).} Let the
adversary corrupt some voters and read the board. If all trustees are
honest, its view can be simulated from the ballots it casts and
$\mathcal{L}(D)$, where $D$ counts the honest voters' rankings and the
rankings in the adversary's accepted ballots (public privacy). If it
also corrupts up to $n - 1$ trustees, its view can be simulated from its
ballots and $D$ (internal privacy).

\emph{Outline.} The argument is the one for
Helios~\cite{bpw12,bcgpw15}, in three steps. (1)~Replace the proofs of
the honest ballots and of the certificate by simulated ones, programming
the oracle at their announcements. \rocq{ballot_zk},
\rocq{ballot3_zk}, and \rocq{tally_zk} say that, for every challenge,
the simulated transcripts have exactly the distribution of the real
ones; the only loss is the chance that a programmed point was queried
before. This is the mechanised part. The certificate simulator takes
the witness paths as input, which is where $\Pi(D)$ enters. (2)~Compute
$D$ without decrypting: take the honest voters' rankings as given, and
extract the ranking of each ballot the adversary casts from its proofs,
which are proofs of knowledge under strong
Fiat--Shamir~\cite{bpw12}. A copy of an honest ballot under another
credential is rejected, because the credential is part of the hash, so
every accepted adversarial ballot carries proofs the adversary made
itself. (3)~Replace the ciphertexts of the honest ballots, and the fresh
ciphertexts of the certificate, by encryptions of $0$. This is
indistinguishable under DDH, and the Pedersen commitments of the batched
ballot are perfectly hiding. After step~(3) the board depends on the
honest rankings only through $w(D)$ and $\Pi(D)$. With corrupted
trustees the adversary learns $D$ in the tally, the honest trustee's
decryption proofs are simulated in the same way, and the board depends
on the honest rankings only through $D$.

Step~(2) is where this outline falls short of a proof. It rewinds once
per adversarial ballot, as the Helios analyses do, and for the batched
ballot it inherits the open step above. We have not written the
reduction out in full, and none of steps~(2) and~(3) is mechanised.
